\subsection{括号}

设 \(r \in N_{K}\)，且 X 是 \(C^{r+1}\) 类流形。

8.5.1. 应用第 8.4 号的定义，取 \(\tau\) 为恒等函子：对每个 Banach 空间 V，有 \(\tau(\mathbf{V}) = \mathbf{V}\)，对 Banach 空间的每个同构 \(\tau(u) = u\)，有 \(u\)。于是 \(\tau(\mathrm{T}(\mathbf{X})) = \mathrm{T}(\mathbf{X})\)。若 \(\xi\) 和 \(\eta\) 是 \(\mathbf{C}^r\) 上的两个 \(\mathbf{X}\) 类向量场，置

\[
[ \xi , \eta ] = \theta_ {\xi}. \eta ;\tag{1}
\]

它是 X 上的 \(C^{r-1}\) 类向量场，称为 \(\xi\) 和 \(\eta\) 的括号。

8.5.2. 映射 \((\xi, \eta) \mapsto [\xi, \eta]\) 是 K-双线性且交错的。若 \(f\) 和 \(g\) 是 \(\mathbf{C}^r\) 上取值于 \(\mathbf{X}\) 的 \(\mathbf{K}\) 类函数，则有：

\[
[ f \xi , g \eta ] = f g [ \xi , \eta ] + (f D _ {\xi} g) \eta - (g D _ {\eta} f) \xi .\tag{2}
\]

8.5.3. 设 F 是 Banach 空间。有

\[
\mathbf {D} _ {[ \xi , \eta ]} = \mathbf {D} _ {\xi} \circ \mathbf {D} _ {\eta} - \mathbf {D} _ {\eta} \circ \mathbf {D} _ {\xi}\tag{3}
\]

在从 \(\mathcal{C}^{r+1}(X;F)\) 到 \(\mathcal{C}^{r-1}(X;F)\) 的映射空间中成立。

若 \(r \geqslant 2\)，且 \(\zeta\) 是 X 上的第三个 \(C^r\) 类向量场，则有

\[
[ [ \xi , \eta ], \zeta ] = [ \xi , [ \eta , \zeta ] ] - [ \eta , [ \xi , \zeta ] ] ^ {(1)}\tag{4}
\]

或者也可写成

\[
[ \xi , [ \eta , \zeta ] ] + [ \eta , [ \zeta , \xi ] ] + [ \zeta , [ \xi , \eta ] ] = 0.
\]

更一般地，若 \(r \geqslant 2\) 且 \(\tau\) 是同构向量函子，则有

\[
\theta_ {[ \xi , \eta ]} = \theta_ {\xi} \circ \theta _ {\eta} - \theta _ {\eta} \circ \theta_ {\xi}\tag{5}
\]

在从 \(\mathcal{C}^{r+1}(\mathrm{X};\mathrm{F})\) 到 \(\mathcal{C}^{r-1}(\mathrm{X};\mathrm{F})\) 的映射空间中成立（其中 \(\mathrm{F} = \tau (\mathrm{T}(\mathrm{X}))\)）。

若 \(r \geqslant \infty\)，则 X 上的 \(C^r\) 类向量场关于该括号构成 K-李代数（LIE, I, § 1, no. 2）。

8.5.4. 设 E 是 Banach 空间，U 是 E 的开子集，\(\xi\) 和 \(\eta\) 是 U 上的两个 \(C^{r}\) 类向量场，并将它们与 \(\mathcal{C}^{r}(U;E)\) 中的元素等同。它们的括号由下式给出：

\[
[ \xi , \eta ] = \mathbf {D} _ {\xi} \eta - \mathbf {D} _ {\eta} \xi .\tag{6}
\]

8.5.5. 设 \((u^1, \ldots, u^n)\) 是 X 在开集 U 上的坐标系。设 \(\xi = \sum a^i \frac{\partial}{\partial u^i}\) 和 \(\eta = \sum b^i \frac{\partial}{\partial u^i}\) 是 U 上的两个 \(C^r\) 类向量场。置 \([\xi, \eta] = \sum c^i \frac{\partial}{\partial u^i}\)。有

\[
c ^ {i} = \sum_ {1 \leqslant j \leqslant n} \left(a ^ {j} \frac {\partial b ^ {i}}{\partial u ^ {j}} - b ^ {j} \frac {\partial a ^ {i}}{\partial u ^ {j}}\right).\tag{7}
\]

8.5.6. 设 \(\varphi: X \to Y\) 是 \(C^{r+1}\) 类流形的态射，\(\xi_1\) 和 \(\xi_2\) 是 X 上的两个向量场，\(\eta_{1}\) 和 \(\eta_{2}\) 是 Y 上的两个向量场，且它们都属于 \(C^r\) 类。若 \(\xi_{1}\) 和 \(\xi_{2}\) 分别与 \(\varphi\) 和 \(\eta_{1}\) 由 \(\eta_{2}\) 相关，则 \([\xi_{1}, \xi_{2}]\) 与 \(\varphi\) 由 \([\eta_{1}, \eta_{2}]\) 相关。

8.5.7. 设 \(\omega\) 是 X 上取值于 Banach 空间 F、属于 \(p\) 类且次数为 \(\mathbf{C}^r\) 的微分形式，\(\xi_0, \ldots, \xi_p\) 是 X 上的 \(\mathbf{C}^r\) 类向量场。有

\[
(\theta_ {\xi_ {0}}. \omega) (\xi_ {1}, \dots , \xi_ {p}) = D _ {\xi_ {0}} \omega (\xi_ {1}, \dots , \xi_ {p}) - \sum_ {i = 1} ^ {p} \omega (\xi_ {1}, \dots , [ \xi_ {0}, \xi_ {i} ], \dots , \xi_ {p})\tag{8}
\]

并且

\[
\begin{array}{l} d \omega (\xi_ {0}, \dots , \xi_ {p}) = \sum_ {i = 0} ^ {p} (- 1) ^ {i} D _ {\xi_ {i}} \omega (\xi_ {0}, \dots , \hat {\xi} _ {i}, \dots , \xi_ {p}) \\ \qquad + \sum_ {i <   j} (- 1) ^ {i + j} \omega ([ \xi_ {i}, \xi_ {j} ], \xi_ {0}, \dots , \hat {\xi} _ {i}, \dots , \hat {\xi} _ {j}, \dots , \xi_ {p}). \end{array}\tag{9}
\]

若 \(\xi\) 和 \(\eta\) 是 \(\mathbf{C}^r\) 类向量场，则有

\[
\theta_ {\xi} \circ i (\eta) - i (\eta) \circ \theta_ {\xi} = i ([ \xi , \eta ])\tag{10}
\]

在从 \(^{r}\Omega^{p}(\mathbf{X};\mathbf{F})\) 到 \(^{r-1}\Omega^{p-1}(\mathbf{X};\mathbf{F})\) 的映射空间中成立。

当 \(p = 1\) 时，公式 (9) 变为：

\[
d \omega (\xi , \eta) = D _ {\xi} \langle \eta , \omega \rangle - D _ {\eta} \langle \xi , \omega \rangle - \langle [ \xi , \eta ], \omega \rangle .\tag{11}
\]

